\section{Reducing the full growing host to characters of the fixed quadratic base}
\label{sec:fixed-base-character-reduction}

The cubic composita in (FG.3) are not abelian over
$\mathbb Q$, but they are abelian over the fixed field
$K=\mathbb Q(\omega)$.  This permits a full character
comparison over $K$, including the growing $2$-power
part of the host.  The integral comparison has a cost;
the point of the following argument is to bound that
cost uniformly rather than discard it.

\begin{theorem}[AC-K25: full character descent over a fixed base]
\label{AC-K25}
Let $F=F_x$ be the host in (FG.3), let
$H=F(\sqrt[3]{269})$, and put $G=\operatorname{Gal}(H/K)$.
For the source range $h\geq1$ one has
\[
 G\simeq C_{2^h}\times C_2^4\times C_3^2,
 \qquad D:=|G|=144\cdot2^h,
 \qquad N:=\exp(G)=3\cdot2^h.                 \tag{FC.1}
\]
For each coefficient prime $\ell$, choose a completion
$O_\ell$ of $\mathbb Z[\zeta_N]$ at a prime above $\ell$,
and set $b_\ell=\operatorname{rank}_{\mathbb Z_\ell}O_\ell$.
Then $b_\ell\leq\varphi(N)=2^h\leq D$.

Let $\Sigma_\ell$ consist of the complex place of $K$,
the places over $\ell$, and the places ramified in $H/K$.
For $e\geq1$ and $\eta\in\widehat G$, define
\[
 \begin{split}
 A_{\eta,e}&=(O_\ell\otimes_{\mathbb Z_\ell}
                  \mathbb Q_\ell/\mathbb Z_\ell)(\eta)(-e),\\
 \mathcal S_{\ell,\eta}(e)&=
 \ker\left(H^1(G_{K,\Sigma_\ell},A_{\eta,e})
 \longrightarrow
 \bigoplus_{v\mid\ell}H^1(K_v,A_{\eta,e})
 \oplus
 \bigoplus_{v\in\Sigma_{\ell,\rm fin},\ v\nmid\ell}
                         H^1(I_v,A_{\eta,e})\right).
 \end{split}                                             \tag{FC.2}
\]
Here $(\eta)$ means that $G_K$ acts through the displayed
character, and $(-e)$ is the usual cyclotomic Tate twist.
All these groups are finite.  For $\ell\ne2,3$,
\[
 b_\ell\log\#S_\ell(H,e)
       =\sum_{\eta\in\widehat G}
                              \log\#\mathcal S_{\ell,\eta}(e).
                                                            \tag{FC.3}
\]
For $\ell=2,3$ there are explicitly specified finite
coefficients $C_\ell(-e)$ such that
\[
 \left|b_\ell\log\#S_\ell(H,e)
       -\sum_{\eta\in\widehat G}
                              \log\#\mathcal S_{\ell,\eta}(e)\right|
       \leq\mathcal C_{K,\Sigma_\ell}(C_\ell(-e)).      \tag{FC.4}
\]
Uniformly in $e$, these two errors satisfy
\[
 \sum_{\ell=2,3}\mathcal C_{K,\Sigma_\ell}(C_\ell(-e))
      \ll D^3\log D_H+D^4\log(2D).                    \tag{FC.5}
\]
The implied constant is absolute for this family.
Consequently the sum of the normalized errors in (FC.4)
over $1\leq e<t_x=\lfloor x^{401}\rfloor$ is
\[
 O_\varepsilon\!\left(
       x^{409+161\varepsilon/5}(\log x)^2\right)
             =o\!\left(\exp(x^\varepsilon/2)\right).  \tag{FC.6}
\]
This comparison is unconditional and covers every
coefficient prime, every character of the full growing
host, and all the indicated weights on the same host.
It does not bound the character groups on the right.
\end{theorem}
\begin{proof}
\proofstep{1. The actual Galois group.}
The disjoint cyclotomic factors proved in (FG.11) give
\[
 \operatorname{Gal}(F/K)
  \simeq(C_2\times C_{2^h})\times C_6\times C_2^2
  \simeq C_{2^h}\times C_2^4\times C_3.
\]
Write $E=K(\sqrt[3]{269})$.
The degree-three extension $E/K$ has no proper
intermediate field.  If $E\cap F\ne K$, it would follow
that $E\subset F$.  But $F/\mathbb Q$ is abelian and
$E/\mathbb Q$ is the nonabelian $S_3$ splitting field
from AM-H03.  A subextension of an abelian Galois
extension is abelian Galois, a contradiction.
Thus $E\cap F=K$.  The compositum of these Galois
extensions has product Galois group over $K$.
Since $\operatorname{Gal}(E/K)=C_3$, this proves (FC.1).
The coefficient-degree inequality follows from the
degree $\varphi(N)$ of the cyclotomic field; every
local completion has degree at most that degree.

\proofstep{2. The integral Fourier map.}
For the moment fix $\ell$ and write $O=O_\ell$.
Use the left regular lattice $L=O[G]$ and the diagonal
lattice $L'=\bigoplus_{\eta\in\widehat G}O(\eta)$.
The Fourier map
\[
 \mathcal F:L\longrightarrow L',\qquad
 \mathcal F\left(\sum_{g\in G}a_g[g]\right)_\eta
                    =\sum_{g\in G}a_g\eta(g)          \tag{FC.7}
\]
is $G$-equivariant.  Character orthogonality gives its
inverse over $\operatorname{Frac}(O)$:
\[
       a_g=D^{-1}\sum_{\eta\in\widehat G}
                          \mathcal F(a)_\eta\eta(g)^{-1}.
                                                            \tag{FC.8}
\]
Indeed the inner character sum at $g^{-1}g'$ is $D$
if $g=g'$ and zero otherwise; a nonidentity group
element is separated by a character, so multiplication
by that character makes the latter sum vanish.
Thus (FC.7) is injective, and its cokernel
$C_\ell$ is killed by $D$.  As a $\mathbb Z_\ell$-module
it is generated by at most $b_\ell D\leq D^2$ elements.
For $\ell\nmid D$ it vanishes.  At the two other primes
it is killed respectively by
\[
                       2^{h+4}\quad\hbox{and}\quad3^2.
                                                            \tag{FC.9}
\]
These are the promised explicit coefficients: they are
the cokernels of the concrete Fourier matrices (FC.7).
No integral idempotent decomposition is asserted at
$2$ or $3$.

Since both lattices are free over $\mathbb Z_\ell$,
tensoring their short exact sequence with
$\mathbb Q_\ell/\mathbb Z_\ell$ gives
\[
 0\longrightarrow C_\ell(-e)
 \longrightarrow A_L
 \longrightarrow A_{L'}\longrightarrow0,
 \qquad
 A_L=L\otimes_{\mathbb Z_\ell}\mathbb Q_\ell/\mathbb Z_\ell(-e),
                                                            \tag{FC.10}
\]
where $A_{L'}=\bigoplus_\eta A_{\eta,e}$.
For completeness, the first term is
$\operatorname{Tor}_1(C_\ell,\mathbb Q_\ell/\mathbb Z_\ell)
\simeq C_\ell$: apply the exact sequence
$0\to\mathbb Z_\ell\to\mathbb Q_\ell
\to\mathbb Q_\ell/\mathbb Z_\ell\to0$ to the finite
module.  The last tensor product of $C_\ell$ with that
divisible group is zero.  This also checks the direction
of the coefficient map in (FC.10).

\proofstep{3. Global and local strict kernels.}
The regular module $A_L$ is the induced module from
$G_H$ with coefficient
$O\otimes_{\mathbb Z_\ell}\mathbb Q_\ell/\mathbb Z_\ell(-e)$.
Shapiro's cochain isomorphism identifies its global
$H^1$ with that over $H$; the local Mackey decomposition
gives the corresponding identifications at decomposition
and inertia subgroups.  In particular at $v\mid\ell$
the restriction is the sum of restrictions at all
$w\mid v$, while restriction at $I_v$ gives restrictions
at the inertia subgroups above $v$, with any repeated
copies immaterial to the vanishing condition.
These are exactly the cochain identifications used in
ALP\_MACHINE, AM-H06, proof of (H06.2)--(H06.4):
the finite coset decomposition and evaluation maps do
not require a cyclic group or a degree prime to $\ell$.
Naturality with respect to restriction verifies the
localization square before taking its kernel.

There is one further point because $\Sigma_\ell$ includes
ramified primes different from $\ell$.  The induced
coefficient is untwisted over $H$, apart from its Tate
twist, and is unramified at these primes.  A global
class over $H$, unramified outside the primes above
$\Sigma_\ell$, whose restrictions to their inertia
groups vanish, descends uniquely to the group unramified
outside $\ell$.  This follows from inflation--restriction:
the extra inertia groups normally generate the kernel,
and a cocycle vanishing on each has zero restriction
to that normal subgroup, since it acts trivially on
the coefficient.  Conversely an inflated class
has precisely those vanishing restrictions.
The strict condition at $\ell$ is unchanged.
Consequently
\[
 \mathcal S(A_L)=O\otimes_{\mathbb Z_\ell}S_\ell(H,e),
 \quad
 \mathcal S(A_{L'})=\bigoplus_\eta\mathcal S_{\ell,\eta}(e),
 \quad
 \log\#\mathcal S(A_L)=b_\ell\log\#S_\ell(H,e).    \tag{FC.11}
\]
The last group is finite by AC-K10; extension by the
finite free coefficient ring commutes with these
cochains and their kernels.

Apply global and local cohomology to (FC.10).  The
kernel of $\mathcal S(A_L)\to\mathcal S(A_{L'})$ is
a quotient of a subgroup of $H^1(G_{K,\Sigma_\ell},C_\ell(-e))$.
To bound the cokernel, first lift a target strict class
globally.  Its only obstruction lies in global $H^2$
of $C_\ell(-e)$.  A class with vanishing obstruction
has a global lift.  The local restrictions of that lift
are in the images of local $H^1(C_\ell(-e))$ because
the target class is strict.  Modulo the changes allowed
by global $H^1(C_\ell(-e))$, their orders are at most
the product of the orders of those local groups.
More formally, the map from the globally liftable
target strict classes modulo the image of the source
strict kernel injects into the quotient of the local
image product by its global image: its zero class
allows subtracting a global lift of the local defects.
Thus
\[
 \begin{split}
 \#\ker&\leq\#H^1(G_{K,\Sigma_\ell},C_\ell(-e)),\\
 \#\operatorname{coker}&\leq
 \#H^2(G_{K,\Sigma_\ell},C_\ell(-e))
 \prod_{v\mid\ell}\#H^1(K_v,C_\ell(-e))
 \prod_{v\in\Sigma_{\ell,\rm fin},\ v\nmid\ell}
                                      \#H^1(I_v,C_\ell(-e)).
 \end{split}                                             \tag{FC.12}
\]
All the bounding groups are finite by AC-K19.
This proves finiteness of the target and (FC.3)--(FC.4).

\proofstep{4. A trivializing field with controlled size.}
For $\ell=2,3$ put $a=v_\ell(D)$ and
$L_\ell=H(\mu_{\ell^a})$.
This is Galois over $K$, unramified outside
$\Sigma_\ell$, trivializes $C_\ell$, and trivializes
every Tate twist of $C_\ell$ since it is killed by
$\ell^a$.  The coefficient ring has scalar action;
its abstract roots of unity introduce no additional
Galois action into this assertion.
Write $m=\varphi(\ell^a)$.
For $\ell=2$ one has $m=2^{h+3}=D/18$;
for $\ell=3$ one has $m=6\leq D$.
Therefore
\[
 [L_\ell:K]\leq D^2,\qquad n_{L_\ell}\leq2D^2,
 \qquad |\Sigma_{\ell,\rm fin}|\leq12,
 \qquad s_{L_\ell}\leq12D^2.                     \tag{FC.13}
\]
The bound twelve is deliberately coarse: the only
rational primes required here are $2,3,7,q,d,269$,
and a quadratic field has at most two places above each.
The cubic extension introduces no other ramification,
as follows from the discriminant of $X^3-269$ and
from the fact that $K\subset F$.

The tensor-order discriminant inequality proved in
(FG.10), together with the cyclotomic discriminant
formula used in (FG.11), gives
\[
 \begin{split}
 \log D_{L_\ell}
 &\leq m\log D_H+2D\log D_{\mathbb Q(\mu_{\ell^a})}\\
 &\leq D\log D_H+2D^2a\log\ell
 \ll D\log D_H+D^2\log(2D).
 \end{split}                                      \tag{FC.14}
\]
In the middle line we used
$\log D_{\mathbb Q(\mu_{\ell^a})}\leq ma\log\ell$.
Now apply (FG.2) with base $K$, with
$r\leq D^2$, and with (FC.13)--(FC.14).
Since $a\log\ell\leq\log D$, its first terms are
$O(D^3\log D_H+D^4\log(2D))$, and its remaining
terms are $O(D^4\log(2D))$.
There are only two primes.  This proves (FC.5),
uniformly in the weight $e$.

\proofstep{5. Insert the source growth parameters.}
Equations (FG.11) and (KG.3)--(KG.4) give
\[
 D=O_\varepsilon(x^{2+8\varepsilon}),\qquad
 \log D_H=O_\varepsilon\!\left(
                  x^{2+41\varepsilon/5}(\log x)^2\right).
\]
Thus $D^3\log D_H$ is
$O_\varepsilon(x^{8+161\varepsilon/5}(\log x)^2)$,
while $D^4\log(2D)$ has no larger order.
There are fewer than $x^{401}$ weights, and dividing
by $b_\ell\geq1$ can only decrease the error.
This proves (FC.6); its final comparison with the
exponential follows exactly as in (FG.4).
\end{proof}

\begin{corollary}[AC-K25-budget: the remaining fixed-base character size]
\label{AC-K25-budget}
Define the finite nonnegative sum
\[
 \mathcal B_x=
 \sum_{1\leq e<t_x}\sum_\ell\frac1{b_\ell}
              \sum_{\eta\in\widehat G}
                         \log\#\mathcal S_{\ell,\eta}(e).
                                                            \tag{FC.15}
\]
Then
\[
 \left|\sum_{1\leq e<t_x}\sum_\ell\log\#S_\ell(H,e)
                              -\mathcal B_x\right|
 =O_\varepsilon\!\left(
                x^{409+161\varepsilon/5}(\log x)^2\right).
                                                            \tag{FC.16}
\]
In particular an $o(\exp(x^\varepsilon/2))$ bound for
either side's nonnegative main sum is equivalent to
the same bound for the other.  At a fixed positive
budget constant the comparison is valid with any
fixed positive relative slack.
\end{corollary}
\begin{proof}
For each fixed $H,e$, AC-K10 makes $S_\ell(H,e)$
zero for all but finitely many $\ell$.
At $\ell\ne2,3$, (FC.3) then forces every character
group to vanish: their logarithms are nonnegative.
The remaining two coefficients give only finitely
many additional terms.  Thus (FC.15) is a finite sum.
Sum (FC.3)--(FC.4), divide by $b_\ell$, and use
(FC.6) and the triangle inequality.  This proves
(FC.16).  Its error is little-o of the stated
budget, giving both assertions about size and slack.
\end{proof}

\paragraph{Meaning of the reduction and the geometric dictionary.}
The base of every character cohomology group in (FC.15)
is the single arithmetic curve
$\operatorname{Spec}\mathcal O_K[\Sigma_\ell^{-1}]$.
The character, its conductor, the coefficient ring,
and the weight still vary.  Abelianity over $K$ makes
the rank-one elliptic-unit framework relevant to this
family; it does not supply a bound for (FC.15).
The main arithmetic gate now has a fixed-base,
rank-one formulation with a proved polynomial total
integral comparison error.  An actual theta/Selmer
tie, the source window quantifiers, and positive slack
remain required for H04.4.

Geometrically, the regular finite etale coefficient
of the cover $\operatorname{Spec}H\to\operatorname{Spec}K$
splits over the coefficient field into Artin characters.
The finite lattices $C_2,C_3$ record the failure of
that decomposition to be integral.  The cover-to-
representation dictionary is SGA~1, Expos\'e~V,
\S\S4 and 7, already inspected in this run; the
cohomological cost is the proved AC-K19, not an
attribution of a numerical inequality to SGA~1.
The direct sums are finite and the Galois base remains
fixed.  No new archive cote or scan is claimed,
and the access limits of Section~\ref{sec:groth-source}
remain in force.
